\documentclass[paper=a4, fontsize=11pt]{scrartcl} % A4 paper and 11pt font size
\usepackage[utf8]{inputenc}

\usepackage[T1]{fontenc} % Use 8-bit encoding that has 256 glyphs
\usepackage[ngerman]{babel} % English language/hyphenation
\usepackage{amsmath,amsfonts,amsthm} % Math packages

\usepackage{graphicx}

\usepackage{sectsty} % Allows customizing section commands
\allsectionsfont{\centering \normalfont\scshape} % Make all sections centered, the default font and small caps

\usepackage{fancyhdr} % Custom headers and footers
\pagestyle{fancyplain} % Makes all pages in the document conform to the custom headers and footers
\fancyhead{} % No page header - if you want one, create it in the same way as the footers below
\fancyfoot[L]{} % Empty left footer
\fancyfoot[C]{} % Empty center footer
\fancyfoot[R]{\thepage} % Page numbering for right footer
\renewcommand{\headrulewidth}{0pt} % Remove header underlines
\renewcommand{\footrulewidth}{0pt} % Remove footer underlines
\setlength{\headheight}{13.6pt} % Customize the height of the header

\numberwithin{equation}{section} % Number equations within sections (i.e. 1.1, 1.2, 2.1, 2.2 instead of 1, 2, 3, 4)
\numberwithin{figure}{section} % Number figures within sections (i.e. 1.1, 1.2, 2.1, 2.2 instead of 1, 2, 3, 4)
\numberwithin{table}{section} % Number tables within sections (i.e. 1.1, 1.2, 2.1, 2.2 instead of 1, 2, 3, 4)

\setlength\parindent{0pt} % Removes all indentation from paragraphs - comment this line for an assignment with lots of text

\newcommand {\horrule}[1]{\rule{\linewidth}{#1}} % Create horizontal rule command with 1 argument of height

\title {
  \normalfont \normalsize
  \textsc{HAW Hamburg} \\ [25pt] % Your university, school and/or department name(s)
  \textsc{Technik und Technologie von vernetzer Systeme} \\ [15pt]
  \horrule{0.5pt} \\[0.4cm] % Thin top horizontal rule
  \huge LoWPAN Networking im IoT \\ [15pt] % The assignment title
  \small Prof. Dr. Fohl \\
  \horrule{2pt} \\[0.5cm] % Thick bottom horizontal rule
}

\author{Fabien Lapok, Matthias Nitsche}

\date{\normalsize\today}

\begin{document}

\maketitle

\section{Aufgabe 1 - Laboreinbindung von Gateways und Sensorknoten}

Formulas...

\begin{align}
  y' &= 10 - 500 \cdot y + 5000 \cdot x
  y(0) &= 1
\end{align}

Text...

\subsection{Subsektion 1}

Some text...

\includegraphics[width=\linewidth]{imgs/cat.jpg}

Some text...

\begin{align}
  y_{n+1} &=  y_n + h \cdot (10 - 500 \cdot y_n + 5000 \cdot x_n) \\
  x_{n+1} &= x_n + h
\end{align}

\section{Aufgabe 2 - RPL Routing im Labornetz}

Some text...

\includegraphics[width=\linewidth]{imgs/cat.jpg}

\section{Aufgabe 3 - Datenverteilung und Messung}

Some text...

\includegraphics[width=\linewidth]{imgs/cat.jpg}

\end{document}

